\documentclass{article}
\usepackage[utf8]{inputenc}
\usepackage{geometry}
\usepackage{amsmath}
\usepackage{amssymb}
\usepackage{amsfonts}
\usepackage[thmmarks, amsmath, thref]{ntheorem}
\usepackage{times} % Fuente Times New Roman
\geometry{a4paper, left=2cm, right=2cm, top=1cm, bottom=1.5cm}
\setlength{\parskip}{0pt} % Espaciado compacto
\setlength{\parindent}{0em}

% Configuración del entorno "Solución"
\theoremstyle{nonumberplain}
\theoremheaderfont{\itshape}
\theorembodyfont{\upshape}
\theoremseparator{.}
\theoremsymbol{\ensuremath{\square}}
\newtheorem{solucion}{Solución}

\begin{document}

\textbf{Problema:} Pruebe que la función \( f(x, y) = -y^3 + 3x^2y \) es armónica. \\

\begin{solucion}
Una función \( C^2 \) que satisface \( \nabla^2 f = 0 \), es decir:
\[
\frac{\partial^2 f}{\partial x^2} + \frac{\partial^2 f}{\partial y^2} = 0
\]
se llama función armónica. Demostraremos que \( f(x, y) = -y^3 + 3x^2y \) cumple esta condición.

**Cálculo de derivadas parciales:**
\[
\frac{\partial f}{\partial x} = \frac{\partial}{\partial x} \left(-y^3 + 3x^2y\right) = 0 + 6xy = 6xy
\]
\[
\frac{\partial^2 f}{\partial x^2} = \frac{\partial}{\partial x} \left(6xy\right) = 6y
\]
\[
\frac{\partial f}{\partial y} = \frac{\partial}{\partial y} \left(-y^3 + 3x^2y\right) = -3y^2 + 3x^2
\]
\[
\frac{\partial^2 f}{\partial y^2} = \frac{\partial}{\partial y} \left(-3y^2 + 3x^2\right) = -6y
\]

**Verificación de la ecuación de Laplace:**
\[
\nabla^2 f = \frac{\partial^2 f}{\partial x^2} + \frac{\partial^2 f}{\partial y^2} = 6y + (-6y) = 0
\]
Por lo tanto, \( f(x, y) = -y^3 + 3x^2y \) es una función armónica, ya que satisface \( \nabla^2 f = 0 \).

\end{solucion}

\end{document}
